% Author: Dr. Matthias Jung, DL9MJ
% Year: 2026
% Fork 50ohm-fr (21/10/2026), montage inchange. L'autoscale ramene l'unite a
% ~0,5 cm sans reduire les composants (1,4 cm) : sur les branches de 2 unites,
% les resistances se chevauchaient. D'ou le \ctikzset, et la boucle superieure
% montee de 2 a 3 unites pour degager l'etiquette de R5 de celle de R8.
\begin{circuitikz}[american]
    \ctikzset{bipoles/length=0.7cm}
    \draw (0, 1.5) node [ocirc](I1){};
    \draw (0,-1.5) node [ocirc](I2){};
    \draw (I1) to [european resistor, l={$R_1=\qty{200}{\ohm}$}] ++(3,0) coordinate(C1) 
               to [european resistor, l={$R_2=\qty{200}{\ohm}$}] ++(3,0) coordinate(C2)
               to [european resistor, l={$R_8=\qty{200}{\ohm}$},-*] ++(3,0) coordinate(C3)
               to [short] ++(0,-3)
               to [european resistor, l={$R_4=\qty{150}{\ohm}$}] (I2);
    \draw (C1) to [short, *-] ++(0,-2)
               to [european resistor, l ={$R_3=\qty{100}{\ohm}$}] ++(3,0)
               to [european resistor, l_={$R_6=\qty{100}{\ohm}$}, -*] ++(0,2);
    \draw (C2) to [european resistor, l_={$R_5=\qty{100}{\ohm}$}] ++(0,3)
               to [european resistor, l ={$R_7=\qty{100}{\ohm}$}] ++(3,0)
               to [short] (C3);
\end{circuitikz}